\section*{历史注记}
\phantomsection
\addcontentsline{toc}{section}{历史注记}

（注意：罗马数字指标是指本注末尾的参考文献。）

正如我们已经看到的（第 I、II、III 章的历史注记），导致微分方程积分法的问题属于 XVII \(^{th}\) 世纪无穷小演算的奠基者们（特别是笛卡儿与巴罗）最早处理的问题之列。从那时起，微分方程理论就没有停止过考验数学家们的敏锐，并且一直是应用分析学中各种各样方法的一片沃土；它所引发的问题远未全部解决，因而人们对它的兴趣经久不衰，以至于它构成了数学与实验科学之间最持久、最富有成果的接触点之一：后者常常从中获得宝贵的帮助，而作为交换又不断提供新的问题。

在一部关于微分方程的现代完整研究所应包含的众多章节中，我们在这里只想阐述其中最初步的两章，它们分别处理存在性定理与线性方程，并假定变量为实变量。这也是我们把这段简短的历史叙述限于这两个方面的原因。从 XVIII \(^{th}\) 世纪初起，数学家们就确信，n 阶微分方程的“一般”积分依赖于 n 个任意常数，而且“一般地”存在一个且仅一个积分，它在变量的一个给定值 \(x_{0}\) 处使函数及其前 n - 1 阶导数取给定的值：支撑这一信念的是（追溯到牛顿的）借助微分方程本身及其前 n 个系数逐个计算解在点 \(x_{0}\) 处的泰勒展开式系数的过程。但直到柯西，才有人研究这样得到的级数的收敛性，并证明其和是微分方程的一个解；当然，那时更谈不上非解析微分方程。在柯西为证明微分方程积分的存在性而发明的各种方法中，我们所遵循的那种方法（(IV) 与 (IV bis)），稍后由利普希茨加以推广，对于非解析方程与积分逼近的情形特别有价值。

线性微分方程属于最早引起注意的问题之列。莱布尼茨与雅各布·伯努利用两次求积积分了一阶线性方程（(I), v. II, p.~731）；任意右端的常系数任意线性方程的一般积分由欧拉得到（(II), p.~296-354）；达朗贝尔用同样的方式求解了常系数线性方程组。稍后，拉格朗日 (III) 处理了 \(n\) 阶线性方程的一般理论，认识到齐次方程的积分是 \(n\) 个积分常数的线性函数，引入了伴随方程，发现了当已知若干特解时齐次方程的降阶法，以及求解非齐次方程的常数变易法。拉格朗日理论中的模糊之处（特别是关于积分的线性无关性）由柯西澄清，他的阐述 (IV bis) 除矩阵记号与初等因子理论所带来的细节改进外，几乎已成定论。

